% Template LaTeX para Informes de Proyectos de Grado de Computación
% InCo, Facultad de Ingeniería, Universidad de la República
% 06/2023

\documentclass[oneside]{prgrado}



%%%%%%%%%%%%%%%%%%%%%%%%%%%%%%%%%%%%%%%%%%%%%%%%%%%%%%%%%%%%%%%%%%%%%%%%%%%%%%%%%%%%%%%%%%%%%%%%
% Datos del Proyecto:
%%%%%%%%%%%%%%%%%%%%%%%%%%%%%%%%%%%%%%%%%%%%%%%%%%%%%%%%%%%%%%%%%%%%%%%%%%%%%%%%%%%%%%%%%%%%%%%%

% título del proyecto (debe escribirse en minúscula, con excepción de la
% letra inicial de la primera palabra y nombres propios)
\title{Título que va a poner a su proyecto de grado}                   

% autores
\author{Nombre de Autor1, Nombre de Autor2, Nombre de Autor3 y Nombre de Autor4}  

% fecha de la defensa
\date{\today}                                 

% supervisor
\supervisor{Nombre Supervisor}                

% cosupervisor (comentar si no tiene)
\cosupervisor{Nombre Co-Supervisor}

% licencia creative commons del documento
% opciones: by, by-sa, by-nd, by-nc, by-nc-sa y by-nc-nd
% opción por defecto: by-nc-nd versión 4.0
\cclicense{by-nc}{4.0}

%%%%%%%%%%%%%%%%%%%%%%%%%%%%%%%%%%%%%%%%%%%%%%%%%%%%%%%%%%%%%%%%%%%%%%%%%%%%%%%%%%%%%%%%%%%%%%%%


%links con colores 
\hypersetup{
  colorlinks   = true,
  urlcolor     = blue,
  linkcolor    = blue,
  citecolor    = red
}


\begin{document}

%%%%%%%%%%%%%%%%%%%%%%%%%%%%%%%%%%%%%%%%%%%%%%%%%%%%%%%%%%%%%%%%%%%%%%%%%%%%%%%%%%%%%%%%%%%%%%%%
% Parte inicial
%%%%%%%%%%%%%%%%%%%%%%%%%%%%%%%%%%%%%%%%%%%%%%%%%%%%%%%%%%%%%%%%%%%%%%%%%%%%%%%%%%%%%%%%%%%%%%%%

\frontmatter % numeración en romanos, capitulos sin numerar

% carátula
\maketitle


%%%%%%%%%%%%%%%%%%%%%%%%%%%%%%%%%%%%%%%%%%%%%%%%%%%%%%%%%%%%%%%%%%%%%%%%%%%%%%%%%%%%%%%%%%%%%%%%

% agradecimientos
\chapter*{Agradecimientos}

Agradecer, siempre es bueno agradecer.



%%%%%%%%%%%%%%%%%%%%%%%%%%%%%%%%%%%%%%%%%%%%%%%%%%%%%%%%%%%%%%%%%%%%%%%%%%%%%%%%%%%%%%%%%%%%%%%%

% resumen
\chapter*{Resumen}

El resumen (200-500 palabras) debe dar una idea completa de todo el
proyecto, mencionando claramente los formalismos, técnicas, herramientas y lenguajes
utilizados. No debe limitarse a describir el problema abordado, sino que debe describir
la solución del problema, con una evaluación de la misma. No debe incluir referencias
bibliográficas ni referencias a otras partes del informe. Tampoco debe utilizar
acrónimos sin explicar su significado.


\hfill \break
\keywords{Template, Proyectos de Grado, Computación}

%%%%%%%%%%%%%%%%%%%%%%%%%%%%%%%%%%%%%%%%%%%%%%%%%%%%%%%%%%%%%%%%%%%%%%%%%%%%%%%%%%%%%%%%%%%%%%%%

% índice
\tableofcontents
\newpage


%%%%%%%%%%%%%%%%%%%%%%%%%%%%%%%%%%%%%%%%%%%%%%%%%%%%%%%%%%%%%%%%%%%%%%%%%%%%%%%%%%%%%%%%%%%%%%%%
% Parte central
%%%%%%%%%%%%%%%%%%%%%%%%%%%%%%%%%%%%%%%%%%%%%%%%%%%%%%%%%%%%%%%%%%%%%%%%%%%%%%%%%%%%%%%%%%%%%%%%

\mainmatter %numeración en arábicos, numerar capítulos 

% introducción
\chapter{Introducción}

Aquí se motiva el trabajo, se plantea y define el problema, se deja claro cuales son los objetivos (general del proyecto, si correspondiese o si está inmerso en un proyecto de mayor alcance, y los específicos), se plantean los resultados esperados, se establecen resumidamente las conclusiones y se describe la
organización general del documento.

Este y los siguientes capítulos pueden incluir referencias, algunos ejemplos de referencias pueden ser \cite{CitekeyArticle}, \cite{CitekeyBook}, \cite{CitekeyInproceedings}, \cite{CitekeyManual} y \cite{CitekeyMisc}. Se sugiere citar usando el formato APA.

%%%%%%%%%%%%%%%%%%%%%%%%%%%%%%%%%%%%%%%%%%%%%%%%%%%%%%%%%%%%%%%%%%%%%%%%%%%%%%%%%%%%%%%%%%%%%%%%

% antecedentes
\chapter{Revisión de antecedentes} 

Revisión de antecedentes (ya sea productos, procesos, publicaciones, etc., a nivel académico o comercial) en el tema del trabajo. Puede incluir además (si es necesario) una breve introducción a los conceptos necesarios para entender el trabajo.

\section{Primera Sección}

Los capítulos pueden incluir secciones como esta, que además pueden hacer referencia a otras secciones, como por ejemplo a la Sección~\ref{sec:segunda}.


\section{Segunda Sección}\label{sec:segunda}

Esta sección además tiene una subsección.

\subsection{La Subsección}

Esta es la subsección.

%%%%%%%%%%%%%%%%%%%%%%%%%%%%%%%%%%%%%%%%%%%%%%%%%%%%%%%%%%%%%%%%%%%%%%%%%%%%%%%%%%%%%%%%%%%%%%%%

% parte central
\chapter{Parte Central}
La parte central del trabajo refiere a lo que es producción propia o aporte del
proyecto de grado, incluyendo las decisiones tomadas. Por ejemplo, puede incluir
los requerimientos, el análisis y el diseño de la solución. Si el proyecto tiene una
implementación, debe describirse en términos de decisiones tomadas en ese sentido.
Los detalles de programación se dejan para los anexos.

Se pueden incluir figuras y tablas en el documento, las mismas deben estar referenciadas en el texto. Por ejemplo, la Figura~\ref{fig:logos} muestra los logos de Facultad de Ingeniería y de la Universidad de la República, mientras que la Tabla~\ref{table:datos} tiene números aleatorios.

\begin{figure}[h!]
    \centering
    \includegraphics[width=\textwidth]{figs/logo-udelar-fing.png}
    \caption{Logos de FIng y UdelaR}
    \label{fig:logos}
\end{figure}

\begin{table}[h!]
\centering
\begin{tabular}{| c | c | c | c |} 
 \hline
 Col1 & Col2 & Col2 & Col3 \\  
 \hline
 1 & 970 & 67 & 941 \\ 
 2 & 668 & 845 & 141 \\
 3 & 800 & 383 & 464 \\
 4 & 143 & 683 & 502 \\
 \hline
\end{tabular}
\caption{Tabla con datos}
\label{table:datos}
\end{table}


%%%%%%%%%%%%%%%%%%%%%%%%%%%%%%%%%%%%%%%%%%%%%%%%%%%%%%%%%%%%%%%%%%%%%%%%%%%%%%%%%%%%%%%%%%%%%%%%

% experimentación
\chapter{Experimentación}
Puede ser necesario incluir un capítulo de Experimentación, incluyendo las pruebas realizadas (casos de prueba) y los resultados obtenidos con su respectivo análisis, que puede incluir comparaciones. 

%%%%%%%%%%%%%%%%%%%%%%%%%%%%%%%%%%%%%%%%%%%%%%%%%%%%%%%%%%%%%%%%%%%%%%%%%%%%%%%%%%%%%%%%%%%%%%%%

%conclusiones
\chapter{Conclusiones y Trabajo Futuro}

En este capítulo se evalúan los resultados alcanzados y
dificultades encontradas, se establece lo que se planteó hacer y lo que se hizo
realmente, cuales fueron los aportes, se muestran posibles extensiones al trabajo, se
realiza una autocrítica de lo que se hizo y lo que faltó (por problemas de tiempo,
recursos, cómo se puede continuar, qué cosas hacer, prioridades, etc.) y se incluye
información sobre la gestión del proyecto, si aplica

%%%%%%%%%%%%%%%%%%%%%%%%%%%%%%%%%%%%%%%%%%%%%%%%%%%%%%%%%%%%%%%%%%%%%%%%%%%%%%%%%%%%%%%%%%%%%%%%
% Parte final
%%%%%%%%%%%%%%%%%%%%%%%%%%%%%%%%%%%%%%%%%%%%%%%%%%%%%%%%%%%%%%%%%%%%%%%%%%%%%%%%%%%%%%%%%%%%%%%%

{ % inicio backmatter

\backmatter % no numerar capítulos 

% referencias bibliográficas

\newpage
\bibliographystyle{apacite}
\bibliography{referencias}

} % fin backmatter

%%%%%%%%%%%%%%%%%%%%%%%%%%%%%%%%%%%%%%%%%%%%%%%%%%%%%%%%%%%%%%%%%%%%%%%%%%%%%%%%%%%%%%%%%%%%%%%%

% anexos

\begin{appendix}

\chapter{Anexo 1}

Los anexos contienen información adjunta al proyecto pero que no es fundamental
para entender el trabajo. Por ejemplo, determinado material de los antecedentes o la
implementación, que en el cuerpo principal del informe se encuentre resumido, aquí
puede presentarse de forma completa. En caso de proyectos de desarrollo de
software, se debería incluir un manual de usuario.

\section{Sección del Anexo}

Los anexos pueden tener secciones.

\end{appendix}


%%%%%%%%%%%%%%%%%%%%%%%%%%%%%%%%%%%%%%%%%%%%%%%%%%%%%%%%%%%%%%%%%%%%%%%%%%%%%%%%%%%%%%%%%%%%%%%%

\end{document}
